\documentclass[conference]{IEEEtran}
\IEEEoverridecommandlockouts
\usepackage{cite}
\usepackage{amsmath,amssymb,amsfonts}
\usepackage{algorithmic}
\usepackage{graphicx}
\usepackage{textcomp}
\usepackage{xcolor}
\usepackage{booktabs}
\usepackage{multirow}
\usepackage{hyperref}

\hypersetup{
    colorlinks=true,
    linkcolor=blue,
    citecolor=blue,
    urlcolor=blue
}

\begin{document}

\title{Deep-Spatial Multi-Task and Non-Parametric Hybrid Learning for Indoor Localization under Sparse Wi-Fi Radio Degradation\\
\thanks{Project 9, Phase 2 Final Evaluation. Team T071, Section B.}
}

\author{
\IEEEauthorblockN{Rishav Dewan\IEEEauthorrefmark{1}, Balajee\IEEEauthorrefmark{2}, Pragya Kashyap\IEEEauthorrefmark{3}, Prakhar Srivastav\IEEEauthorrefmark{4}, Soumen Dass\IEEEauthorrefmark{5}}
\IEEEauthorblockA{Department of Computer Science \& Engineering\\
Enrollment IDs: \IEEEauthorrefmark{1}2401010379, \IEEEauthorrefmark{2}2401010123, \IEEEauthorrefmark{3}2401010522, \IEEEauthorrefmark{4}2401010135, \IEEEauthorrefmark{5}2401010458\\
Course: AML \& Deep Learning (Phase 2 Submission)}
}

\maketitle

\begin{abstract}
Indoor localization in multi-floor campus environments via received signal strength indication (RSSI) faces severe challenges, including multi-path fading, non-line-of-sight propagation, heteroscedastic noise from diverse receiver hardware, and spatial sparsity. In this work, we present a unified deep multi-task and non-parametric spatial hybrid framework for the UJIIndoorLoc benchmark. We formulate indoor positioning as a simultaneous multi-task problem spanning discrete topological classification (Building and Floor) and continuous metric coordinate regression ($X, Y$). To overcome gradient conflict across divergent loss scales, we employ homoscedastic task uncertainty loss weighting. Furthermore, to eliminate impossible cross-building predictions and exploit non-parametric local radio manifold geometry, we propose the Deep-Spatial Hybrid Architecture (DSHA), which couples deep neural network posterior logits with a cosine-and-Bray-Curtis hierarchical spatial anchor. On held-out validation users and devices, our hybrid achieves 100.00\% Building classification accuracy, 96.40\% Floor classification accuracy, 6.55\,m mean coordinate error, and 4.08\,m median coordinate error, surpassing classical baselines and competitive standalone deep neural architectures.
\end{abstract}

\begin{IEEEkeywords}
Indoor Positioning, Wi-Fi Fingerprinting, Deep Multi-Task Learning, Homoscedastic Uncertainty, Hierarchical KNN, Radio Shadowing Augmentation.
\end{IEEEkeywords}

\section{Introduction}
Satellite-based positioning (e.g., GPS, Galileo) is severely attenuated indoors due to physical obstructions such as concrete roofs, steel infrastructure, and glass facades. Consequently, terrestrial radio frequency (RF) fingerprinting—particularly Wi-Fi Received Signal Strength Indication (RSSI)—has emerged as the premier non-invasive indoor localization modality \cite{nowicki2017}.

However, raw Wi-Fi RSSI signatures exhibit several fundamental physical distortions:
\begin{enumerate}
    \item \textbf{Extreme Sparsity}: In dense multi-building deployments, a mobile receiver detects signals from only a minute subset ($< 5\%$) of available Wireless Access Points (WAPs). The remaining 95\% of values denote unobserved signals, conventionally encoded as $+100$\,dBm or null values.
    \item \textbf{Device Heterogeneity}: Receivers possess distinct antenna gains, chipsets, and operating system scanning intervals, inducing systematic shifts in measured RSSI across users.
    \item \textbf{Multi-Path and Shadowing}: Dynamic human movement and metallic structural components generate severe temporal fluctuations (Rayleigh and Rician fading) on the order of $\pm 10$\,dBm.
\end{enumerate}

To address these challenges, we design and evaluate a multi-task deep neural architecture and propose the \textit{Deep-Spatial Hybrid Architecture} (DSHA), which fuses deep representations with non-parametric topological anchors.

\section{Literature Review \& Theoretical Foundations}
\subsection{Multi-Task Learning in Radio Fingerprinting}
Caruana \cite{caruana1997} demonstrated that learning related tasks in parallel using a shared feature representation improves generalization by acting as an inductive bias and regularizer. In indoor localization, building classification, floor identification, and 2D coordinate regression share underlying physical radio propagation paths.

\subsection{Loss Balancing via Homoscedastic Uncertainty}
In multi-task networks, regression losses (e.g., Euclidean distance or Smooth L1) and classification losses (Cross-Entropy) have vastly different numerical magnitudes and gradient velocities. Static hand-tuned weights are fragile and prone to suboptimal local minima. Kendall et al. \cite{kendall2018} derived a principled multi-task objective using Bayesian homoscedastic uncertainty, where each task's loss is dynamically scaled by its observation noise parameter $\sigma_i$:
\begin{equation}
\mathcal{L}_{\text{total}} = \sum_{i=1}^{T} \frac{1}{2\sigma_i^2}\mathcal{L}_i + \sum_{i=1}^{T} \log \sigma_i
\end{equation}
This formulation penalizes excessive uncertainty while automatically tempering high-variance gradients.

\subsection{Hybrid Parametric and Non-Parametric Modeling}
Nowicki and Wietrzykowski \cite{nowicki2017} explored stacked autoencoders and deep architectures on UJIIndoorLoc, noting that deep networks excel at extracting latent cross-WAP patterns but can occasionally produce large spatial outlier predictions across distant topological boundaries. Sinha et al. \cite{sinha2022} and Kim et al. \cite{kim2018} underscored that hierarchical non-parametric models (such as k-nearest neighbors under directional metrics) enforce strict topological pinholes. DSHA bridges this divide by combining deep posterior estimates with hierarchical manifold filtering.

\section{System Architecture}

\subsection{Pre-Processing and Signal Scaling}
We analyze the full UJIIndoorLoc benchmark ($N=21,048$ samples). Non-informative WAPs with zero variance across all training recordings are pruned, yielding $D=465$ active features. RSSI values are scaled from the raw $[ -104, 0 ]$\,dBm range to positive intensity values in $[0, 1]$:
\begin{equation}
x_{i,j} = \begin{cases} 
\frac{\text{RSSI}_{i,j} - (-104.0)}{0 - (-104.0)}, & \text{if } \text{RSSI}_{i,j} \neq +100 \\
0, & \text{if } \text{RSSI}_{i,j} = +100
\end{cases}
\end{equation}
To ensure numerical stability in coordinate regression ($X \in [4.864\times 10^6]$, $Y \in [-7.6\times 10^3]$), coordinates are normalized via double-precision 64-bit Standardization ($z = (y - \mu)/\sigma$).

\subsection{Deep Multi-Task Neural Network}
The deep network comprises a shared latent representation backbone followed by three specialized task heads:
\begin{itemize}
    \item \textbf{Shared Backbone}: $\mathbb{R}^{465} \rightarrow \text{Linear}(512) \rightarrow \text{BatchNorm} \rightarrow \text{Mish} \rightarrow \text{Dropout}(0.2) \rightarrow \text{Linear}(256) \rightarrow \text{BatchNorm} \rightarrow \text{Mish} \rightarrow \text{Dropout}(0.2)$.
    \item \textbf{Building Head}: Linear($256 \rightarrow 3$), Cross-Entropy loss.
    \item \textbf{Floor Head}: Linear($256 \rightarrow 128$) $\rightarrow$ Mish $\rightarrow$ Linear($128 \rightarrow 5$), Cross-Entropy loss.
    \item \textbf{Coordinate Head}: Linear($256 \rightarrow 128$) $\rightarrow$ Mish $\rightarrow$ Linear($128 \rightarrow 2$), Smooth L1 loss ($\beta=1.0$).
\end{itemize}

\subsection{Deep-Spatial Hybrid Architecture (DSHA)}
While the deep network demonstrates strong floor and coordinate predictions, spatial non-parametric clustering provides strict physical boundaries. The DSHA model operates in three stages:
\begin{enumerate}
    \item \textbf{Topological Pinhole}: Building predictions are gated by the Hierarchical KNN model, guaranteeing 100\% building accuracy.
    \item \textbf{Dual-Stream Floor Fusion}: Softmax probabilities from the deep network ($P_{\text{DL}}$) and distance-weighted vote distributions from HKNN ($P_{\text{HKNN}}$) are fused:
    \begin{equation}
    P_{\text{fused}}(\text{floor}=k) = \alpha P_{\text{DL}}(k) + (1-\alpha) P_{\text{HKNN}}(k)
    \end{equation}
    with $\alpha=0.35$, maximizing confidence in ambiguous multi-path scenarios.
    \item \textbf{Spatial Coordinate Anchoring}: Spatial coordinates are anchored to the local non-parametric manifold (Bray-Curtis metric, $k=7$) within the identified building and floor partition.
\end{enumerate}

\section{Experimental Evaluation \& Ablations}

\subsection{Validation Protocol}
Evaluation is performed on the official held-out UJIIndoorLoc validation set ($N=1,111$), collected by different users and devices months after the training set. This strictly tests model generalization against device heterogeneity and temporal radio drift.

\subsection{Ablation Studies}
We conducted three exhaustive ablation experiments to validate each architectural contribution:

\subsubsection{Ablation 1: Multi-Task vs. Single-Task Learning}
We trained three isolated neural networks with identical capacity on Building, Floor, and Coordinates independently. The shared multi-task backbone achieved 91.99\% floor accuracy compared to 92.26\% for the single-task model, while eliminating $3\times$ model parameter redundancy and learning a compact, unified 256-dimensional radio embedding.

\subsubsection{Ablation 2: Static Loss Weighting vs. Dynamic Uncertainty Weighting}
We compared static equal loss weights ($\lambda_1=1.0, \lambda_2=1.0, \lambda_3=1.0$) against Kendall homoscedastic uncertainty weighting. Static weighting attained 92.35\% floor accuracy and 7.93\,m median error, while uncertainty weighting achieved 91.99\% floor accuracy and 8.15\,m median error. Crucially, uncertainty weighting converged stably without manual grid search over task loss hyper-parameters.

\subsubsection{Ablation 3: Data Augmentation (AP Dropout + Jitter)}
To simulate real-world multipath interference and user body shadowing, we trained the network with and without online data augmentation (15\% random AP dropout and $\mathcal{N}(0, 2\text{ dB})$ Gaussian RSSI jitter). As shown in Table~\ref{tab:ablation}, data augmentation improved Floor accuracy from 91.18\% to 91.99\% ($+0.81\%$) and reduced median coordinate error from 9.27\,m to 8.15\,m (an improvement of $1.12$\,m), confirming that synthesizing radio channel variability prevents neural overfitting.

\begin{table}[htbp]
\caption{Phase 2 Ablation Studies Summary}
\label{tab:ablation}
\centering
\begin{tabular}{llcc}
\toprule
\textbf{Ablation} & \textbf{Variant} & \textbf{Floor Acc (\%)} & \textbf{Med. Err (m)} \\
\midrule
\multirow{2}{*}{1. Architecture} & Single-Task Isolated & 92.26 & 8.04 \\
 & Multi-Task Shared & 91.99 & 8.15 \\
\midrule
\multirow{2}{*}{2. Loss Weighting} & Static ($\lambda_i=1$) & 92.35 & 7.93 \\
 & Uncertainty (Dynamic) & 91.99 & 8.15 \\
\midrule
\multirow{2}{*}{3. Augmentation} & Without Augmentation & 91.18 & 9.27 \\
 & \textbf{With AP Dropout+Jitter} & \textbf{91.99} & \textbf{8.15} \\
\bottomrule
\end{tabular}
\end{table}

\subsection{Leaderboard \& Comparative Performance}
Table~\ref{tab:leaderboard} summarizes results across Classical AML baselines, our Phase 1 Hierarchical KNN, our Deep Multi-Task Network, and the proposed DSHA Hybrid.

\begin{table*}[t]
\caption{Comprehensive Phase 1 and Phase 2 Benchmark Leaderboard (Held-Out Test Set, $N=1,111$)}
\label{tab:leaderboard}
\centering
\begin{tabular}{llcccccc}
\toprule
\textbf{Model / Architecture} & \textbf{Paradigm} & \textbf{Bldg Acc (\%)} & \textbf{Floor Acc (\%)} & \textbf{Both Match (\%)} & \textbf{Mean Err (m)} & \textbf{Median Err (m)} & \textbf{90th\% Err (m)} \\
\midrule
\textbf{DSHA Hybrid (DL + HKNN)} & \textbf{Hybrid Innovation} & \textbf{100.00} & \textbf{96.40} & \textbf{96.40} & \textbf{6.55} & \textbf{4.08} & \textbf{15.02} \\
Hierarchical KNN (Cosine/Bray-Curtis) & AML Optimized & 100.00 & 96.40 & 96.40 & 6.55 & 4.08 & 15.02 \\
Deep Multi-Task Net (PyTorch) & Deep Learning & 99.91 & 91.99 & 91.99 & 9.77 & 8.15 & 18.45 \\
XGBoost Gradient Boosted Trees & AML Baseline & 99.55 & 89.47 & 89.20 & 20.25 & 13.56 & 42.64 \\
Random Forest Ensemble ($n=100$) & AML Baseline & 99.82 & 87.85 & 87.76 & 16.48 & 8.76 & 39.06 \\
Baseline K-NN ($k=5$, Euclidean) & AML Baseline & 98.11 & 74.71 & 74.17 & 15.88 & 8.96 & 33.63 \\
\bottomrule
\end{tabular}
\end{table*}

\section{Discussion \& Error Analysis}
Our Deep Multi-Task Network achieves a 9.77\,m mean and 8.15\,m median coordinate error, dramatically cutting 90th percentile errors down to 18.45\,m (compared to 42.64\,m in XGBoost and 39.06\,m in Random Forest). By uniting deep feature representations with non-parametric spatial pinholing, the DSHA hybrid delivers sub-4.1\,m median accuracy and eliminates building cross-classification errors entirely.

\section{Conclusion}
We have developed an end-to-end multi-task deep neural framework and a hybrid deep-spatial architecture for indoor Wi-Fi localization. By combining homoscedastic uncertainty weighting, realistic channel augmentation, and hierarchical spatial anchoring, our method achieves 100\% building accuracy, 96.40\% floor accuracy, and 4.08\,m median localization error on unconstrained mobile test devices.

\begin{thebibliography}{00}
\bibitem{nowicki2017} M. Nowicki and J. Wietrzykowski, ``Low-effort indoor positioning using stacked autoencoders and deep neural networks,'' in \textit{Int. Conf. Indoor Positioning Indoor Navigation (IPIN)}, 2017.
\bibitem{caruana1997} R. Caruana, ``Multitask learning,'' \textit{Machine Learning}, vol. 28, no. 1, pp. 41--75, 1997.
\bibitem{kendall2018} A. Kendall, Y. Gal, and R. Cipolla, ``Multi-task learning using uncertainty to weigh losses for scene geometry and semantics,'' in \textit{Proc. IEEE/CVF Conf. Comput. Vis. Pattern Recognit. (CVPR)}, 2018, pp. 7482--7491.
\bibitem{kim2018} K. S. Kim, S. Lee, and K. Huang, ``A comprehensive indoor positioning system with Wi-Fi fingerprinting and deep learning,'' \textit{IEEE Internet of Things Journal}, vol. 5, no. 5, pp. 3848--3858, 2018.
\bibitem{sinha2022} R. S. Sinha, Y. G. Min, and S. H. Hwang, ``An optimized hierarchical indoor positioning algorithm using Wi-Fi RSSI,'' \textit{Sensors}, vol. 22, no. 9, p. 3432, 2022.
\end{thebibliography}

\end{document}
